\documentclass{article}

%%% ready for submission (anonymized, double-blind) %%%
\usepackage{cpal_2027}

%%% to compile a preprint version, e.g., for submission to arXiv, add the [preprint] option %%%
%\usepackage[preprint]{cpal_2027}

%%% to compile a camera-ready version, add the [final] option %%%
%\usepackage[final]{cpal_2027}

\usepackage{url}
\usepackage{amsmath,amssymb,amsthm}
\usepackage{graphicx}
\usepackage{booktabs}
\usepackage{multirow}
\usepackage{xcolor}
\usepackage{pgfplots}
\pgfplotsset{compat=1.18}
\usetikzlibrary{patterns}

\newtheorem{proposition}{Proposition}
\newtheorem{observation}{Observation}

\newcommand{\RR}{\mathbb{R}}
\newcommand{\theorig}{\theta^\star}
\newcommand{\gvec}{g}
\newcommand{\Hmat}{H}

\title{Direction, Not Magnitude: Why Different Compression Mechanisms\\ Harm Different Groups}

\author{%
  Anonymous Author(s)\\
  Affiliation\\
  \texttt{email@domain.com}
}

\begin{document}

\maketitle

\begin{abstract}
Model compression trades a small accuracy cost for efficiency, but recent work shows this cost falls unevenly across demographic groups, and attributes the resulting harm to how much a group's loss landscape is perturbed: groups with larger gradient norms or higher-curvature loss surfaces absorb more excess loss. We show that this norm-based view is incomplete. Because it depends only on perturbation \emph{magnitude}, it predicts the identical ranking of harmed groups regardless of which compression mechanism produces the perturbation, and we show this prediction is false. We use exact Hessian-vector products to compute the full directional second-order term, calibrating unstructured pruning, structured (channel) pruning, low-rank truncation, and quantization to the same total weight-perturbation magnitude. We find that these mechanisms disagree on which group suffers most in every trial. Any predictor sensitive to perturbation \emph{direction}, whether it uses the full Hessian, only its diagonal, or only the bare first-order term $g_a^\top\delta$, correctly identifies the most-harmed group in $83.3\%$ of these disagreement cases, against $25.0\%$ for two predictors built only from perturbation \emph{magnitude} (an isotropic trace-based score matching prior quantization-fairness theory, and the norm bound itself). Curvature, whether exact or diagonal-only, adds no measurable predictive power here beyond simply getting the sign of the first-order term right: we find the three direction-aware predictors statistically indistinguishable from one another, so the decisive distinction is between seeing $\delta$'s direction at all and not, not between how much directional detail a predictor retains. We further show this theory is itself regime-dependent. It is accurate only when the weight perturbation is small relative to the model, and it degrades exactly in the high-sparsity regime where compression-fairness harms are largest in practice. Our results reframe the choice of \emph{compression mechanism}, not only compression \emph{severity}, as a fairness-relevant design decision. We show that even the cheapest such diagnostic, a single gradient inner product requiring no Hessian-vector product at all, anticipates it as well as the curvature-aware predictors, within the regime where the underlying theory is valid.
\end{abstract}

\section{Introduction}
\label{sec:intro}

Network compression (pruning, quantization, low-rank factorization) is the standard route to deploying large models under memory, latency, or energy constraints. A substantial literature has shown that the resulting small average accuracy loss is not distributed evenly: compression disproportionately harms performance on minority subpopulations and rare or atypical examples \citep{hooker2020characterising,tran2022pruning,iofinova2023bias}. The leading explanation for \emph{why} comes from a second-order Taylor argument: a group's excess loss under compression is bounded by a term involving the norm of that group's loss gradient and the top eigenvalue of its loss Hessian, both evaluated at the uncompressed optimum \citep{tran2022pruning}. Groups with larger gradients or more curvature, typically small or atypical groups sitting closer to the decision boundary, absorb more harm. The same style of argument, using the \emph{trace} of the group Hessian rather than its top eigenvalue, has since been used to explain disparate impact under quantization specifically \citep{liu2025understanding,bellam2024causes,bellam2025explaining}.

These theories share a structural property that has gone unexamined: they bound excess loss using only the \emph{magnitude} of the weight perturbation, $\|\delta\|$, where $\delta = \theta_{\text{compressed}} - \theorig$. A bound of the form $\|g_a\| \cdot \|\delta\| + \tfrac{1}{2}\lambda(H_a)\|\delta\|^2$ is agnostic to the \emph{direction} of $\delta$: it would assign the same ranking of harmed groups to any two compression mechanisms that happen to produce perturbations of the same size, no matter how differently those perturbations are shaped in weight space. Yet different compression mechanisms are not interchangeable in this way. Unstructured magnitude pruning zeroes individually small-magnitude weights; structured pruning removes entire correlated channels; low-rank truncation discards an entire tail subspace of a weight matrix. These are geometrically distinct operations on the same weight space, and there is no a priori reason they should harm the same group most.

\textbf{This paper shows they do not.} We revisit the second-order expansion of a group's excess loss without discarding directional information: the exact term is $g_a^\top \delta + \tfrac{1}{2}\delta^\top H_a \delta$, not its norm-based upper bound. We show that at a stationary point of the (population-weighted) training loss, the norm-based bound's group ranking is \emph{mechanism-independent by construction}: it can only ever reproduce a single, fixed ordering of groups, set once at $\theorig$ and carried unchanged into every compression mechanism. Only the directional term can predict a ranking that \emph{changes} with the mechanism. We treat this as a falsifiable, nearly zero-cost empirical test: compression mechanisms that are calibrated to displace the model by the same total amount should, under the norm-based theory, always harm the same group most; under the directional theory, they need not.

We test this on a controlled benchmark with a known spurious correlation, comparing unstructured pruning, structured channel pruning, low-rank (SVD) truncation, and quantization at matched relative weight-perturbation norm. We find that which group is harmed most genuinely changes across mechanisms, in every trial. Three direction-aware predictors anticipate this correctly in $83.3\%$ of these disagreement cases: the exact second-order term, its diagonal-only approximation used in prior mitigation methods, and even the bare first-order term $g_a^\top\delta$ alone. Two predictors built only from perturbation \emph{magnitude} (an isotropic/trace-based score matching prior quantization-fairness theory, and the norm bound itself) manage only $25.0\%$. We find the three direction-aware predictors statistically indistinguishable from each other: curvature, exact or diagonal-only, adds no measurable power here beyond getting the \emph{sign} of the first-order term right. The decisive contrast is between seeing $\delta$'s direction at all and not seeing it, not between \emph{how much} directional detail a predictor retains. This is good news practically: a single gradient per group, not a Hessian-vector product, already delivers this predictor's benefit (Section~\ref{sec:discussion}). We also document a calibration pitfall relevant to any similar comparison: matching mechanisms on \emph{overall accuracy loss} rather than on weight-perturbation magnitude silently permits perturbations large enough to leave the regime where any second-order theory, including ours, is valid. We characterize this failure mode directly.

\textbf{Contributions.}
\begin{itemize}
  \item We show the existing norm-based disparate-impact theory, read as a predictor of which group a compression mechanism harms most, is mechanism-agnostic by construction, and give the full (un-bounded) second-order term that is not.
  \item We show the isotropic/trace-based account used for quantization-fairness \citep{liu2025understanding,bellam2024causes,bellam2025explaining} shares this same mechanism-blindness at matched perturbation magnitude, since it is also a function of $\delta$ only through $\|\delta\|$; a Gauss--Newton analysis motivates, without proving for the exact Hessian, a further regime (confidently misclassified examples) where this approximation may be least reliable.
  \item We empirically compare five predictors (the norm bound, the isotropic/trace-based score, a gradient-retaining diagonal-curvature score, a first-order-only score, and the full directional predictor) on a controlled spurious-correlation benchmark, across four compression mechanisms (unstructured pruning, channel pruning, low-rank truncation, and quantization) calibrated to the same relative weight-perturbation norm and three random seeds. We find that the three direction-aware predictors (full, diagonal, first-order-only) are statistically indistinguishable from one another while both magnitude-only predictors perform far worse: the decisive factor is seeing $\delta$'s direction at all, not how much curvature information a direction-aware predictor retains.
  \item We identify, and correct, a calibration pitfall relevant to any similar cross-mechanism comparison: matching mechanisms by their effect on \emph{overall} loss, rather than by weight-perturbation norm directly, can silently select perturbations large enough to leave the regime where any second-order account is valid.
\end{itemize}

\section{Related Work}
\label{sec:related}

\textbf{Disparate impact of compression.} An empirical literature documents that pruning, quantization, and distillation disproportionately harm minority subpopulations and rare classes \citep{hooker2020characterising,iofinova2023bias}, with evidence extending to large language models under more diverse compression methods, including a documented case where the pruning method that best preserves perplexity amplifies bias the most \citep{wu2026asymmetric,rath2026weight}. \citet{tran2022pruning} give the first mechanistic account for pruning specifically, bounding group excess loss via the group gradient norm and the \emph{maximum eigenvalue} of the group Hessian, a worst-case curvature direction that, notably, does not depend on the actual perturbation direction, and propose a Lagrangian mitigation that constrains per-group loss. We build directly on this bound. \citet{liu2025understanding} and \citet{bellam2024causes,bellam2025explaining} give an analogous account for quantization, attributing harm to group gradient norm and Hessian \emph{trace}, which implicitly assumes an isotropic (noise-like) perturbation. We show in Section~\ref{sec:isotropic-special-case} that this is the isotropic special case of the more general directional theory we propose, and in Section~\ref{sec:results} that it is a substantially weaker predictor once perturbation direction is allowed to vary, as it does across structured compression mechanisms.

\textbf{Classical pruning-saliency criteria.} Independently of any fairness motivation, four decades of pruning work have asked which Taylor term of the (pooled, ungrouped) loss best predicts the cost of removing a weight or filter. The answer has moved back and forth between first-order and second-order, and between diagonal and full-Hessian, criteria. \citet{lecun1990obd} rank weights by a diagonal-Hessian saliency, explicitly assuming off-diagonal curvature is negligible. \citet{hassibi1993obs} show this assumption can mislead and rank by the full inverse Hessian instead, at substantially higher cost. \citet{singh2020woodfisher} revisit that trade-off with a Fisher-information approximation to the full, not diagonal, Hessian, making it tractable at modern scale. A separate line finds that a first-order criterion, a product of gradient and weight with no curvature term at all, is enough to rank connections for single-shot pruning at initialization \citep{lee2019snip} or filters for iterative pruning after training \citep{molchanov2019importance}, matching or beating the heavier curvature-based scores each paper compares against. We find that a bare first-order term matches the full and diagonal predictors as well (Section~\ref{sec:baseline-comparison}). This is the same pattern this literature has reported for \emph{pooled} loss, and we show it holds for \emph{group-disaggregated} loss too, in exactly the setting where curvature seemed most likely to matter: predicting which population a given mechanism harms, not only how much the population loss moves.

\textbf{Mitigation via group-conditioned saliency.} A line of work on fair dermatological diagnosis uses a per-group, per-weight saliency score to decide \emph{which} weights to compress, for pruning \citep{wu2022fairprune}, channel pruning specifically \citep{kong2024channel}, quantization \citep{guo2024fairquantize}, and low-rank factorization \citep{guo2025fairlrf}; all four use a diagonal (per-coordinate) Hessian approximation in the OBD style above, dropping the gradient term and off-diagonal curvature. \citet{proskurina2025fairgptq} instead augments GPTQ's \citep{frantar2023gptq} layer-local quantization objective with a bias-aware term built from activation differences between contrastive sentence pairs: a full but layer-local, input-space Hessian; it is specific to quantization and is not a group-loss saliency score at all. \citet{hashemizadeh2024balancing} constrains disparate impact during sparse training via constrained optimization rather than a saliency score. These methods establish that group-conditioned second-order information is a useful \emph{mitigation} tool. Our empirical comparison (Section~\ref{sec:results}) includes a gradient-retaining diagonal variant of this family of scores, rather than the exact zero-gradient form these methods use, since our question is whether perturbation \emph{direction} beyond a scalar magnitude is informative at all, not whether to match a specific prior method's exact formula. \citet{guo2026faircompressagent} is, to our knowledge, the closest existing work to a cross-mechanism treatment: an agentic search procedure that reuses a diagonal saliency score across pruning, quantization, and low-rank factorization to configure a compressed model. It treats these mechanisms as interchangeable operators under one shared score and selects among them by measured search rather than prediction, reporting composition effects (e.g.\ that adding quantization to pruning can restore feasibility where either alone fails) that it leaves unexplained.

\section{Background: The Norm-Based Disparate-Impact Bound}
\label{sec:background}

Let $\theorig$ be the parameters of a trained model at a stationary point of the (population-weighted) training loss $L(\theta) = \sum_a w_a L_a(\theta)$, where $L_a$ is the loss restricted to group $a$ and $w_a$ its population weight, $\sum_a w_a = 1$. Write $g_a = \nabla_\theta L_a(\theorig)$ and $H_a = \nabla^2_\theta L_a(\theorig)$ for the group's gradient and Hessian at $\theorig$. Compressing the model produces new parameters $\theta_{\text{compressed}} = \theorig + \delta$. A second-order Taylor expansion gives group $a$'s \emph{excess loss}
\[
R(a) \;=\; L_a(\theorig + \delta) - L_a(\theorig) \;=\; g_a^\top \delta + \tfrac{1}{2}\delta^\top H_a \delta + O(\|\delta\|^3).
\]
\citet{tran2022pruning} bound this via Cauchy--Schwarz and the top eigenvalue $\lambda(H_a)$ of $H_a$:
\begin{equation}
R(a) \;\le\; \|g_a\| \cdot \|\delta\| + \tfrac{1}{2}\lambda(H_a) \cdot \|\delta\|^2 + O(\|\delta\|^3). \label{eq:norm-bound}
\end{equation}
This bound is tight for a specific $\delta$ (the top eigenvector of $H_a$, aligned with $g_a$) but is an inequality for any other $\delta$, and in particular discards the sign and direction of $g_a^\top\delta$ and the off-diagonal structure of $\delta^\top H_a \delta$.

\begin{observation}
\label{obs:mechanism-agnostic}
The right-hand side of Eq.~\eqref{eq:norm-bound} depends on the compression mechanism \emph{only} through the scalar $\|\delta\|$. For two mechanisms producing perturbations $\delta_1, \delta_2$ with $\|\delta_1\| = \|\delta_2\|$, the norm-based bound assigns \emph{identical} excess-loss bounds to every group under both mechanisms, and therefore predicts the same group ranking regardless of mechanism.
\end{observation}

This is immediate from Eq.~\eqref{eq:norm-bound}. \citet{tran2022pruning} present it as an upper bound, not as a claim about which group's \emph{actual} excess loss is largest; but the bound is the only mechanism-level quantity it offers, and mitigation practice built on it (Section~\ref{sec:related}) necessarily uses it as a ranking. Read that way, as a predictor of which group a given compression mechanism harms most, its consequence is easy to miss: any two mechanisms matched on total perturbation size are, under this reading, interchangeable in terms of \emph{who} they harm most. Section~\ref{sec:theory} shows the exact (non-bounded) second-order term need not share this property, and Section~\ref{sec:results} tests this directly.

\section{Theory: Directional Alignment Determines Mechanism-Specific Harm}
\label{sec:theory}

\subsection{The full second-order term, and why it can flip rankings}

We define the (un-bounded) second-order excess loss $\widehat{R}(a) = g_a^\top \delta + \tfrac12 \delta^\top H_a \delta$: the same Taylor term as Eq.~\eqref{eq:norm-bound}, dropping only the identical $O(\|\delta\|^3)$ remainder, but kept as a signed value rather than bounded via Cauchy--Schwarz. Unlike the norm bound, $\widehat{R}(a)$ depends on $\delta$'s full direction, not only its magnitude. Different compression mechanisms produce perturbations $\delta$ with qualitatively different directional structure for the same target $\|\delta\|$:
\begin{itemize}
  \item \textbf{Unstructured magnitude pruning} sets $\delta_i = -\theorig_i$ on the coordinates $i$ with smallest $|\theorig_i|$: a perturbation supported on a sparse, magnitude-selected coordinate subset.
  \item \textbf{Structured (channel) pruning} sets $\delta$ to zero whole output channels (rows of a weight tensor) at once: a perturbation with strong coordinate \emph{correlation} within each removed channel.
  \item \textbf{Low-rank truncation} sets $\delta$ to lie in the tail (small-singular-value) subspace of a weight matrix, viewed as a linear map: a perturbation shaped by the matrix's spectral geometry rather than individual weight magnitudes.
\end{itemize}
Because $g_a$ and $H_a$ are themselves fixed, group-specific, mechanism-independent objects (evaluated once, at $\theorig$), whether a given mechanism harms group $a$ more than group $b$ is decided by how well each mechanism's particular $\delta$ \emph{aligns} with $g_a$ and the quadratic form of $H_a$, relative to its alignment with $g_b$ and $H_b$. Observation~\ref{obs:mechanism-agnostic} shows this alignment information is exactly what Eq.~\eqref{eq:norm-bound} throws away.

\begin{proposition}[Stationarity constrains, but does not fix, direction]
\label{prop:stationarity}
If $\theorig$ is a stationary point of the \emph{unregularized, population-weighted training loss} $L$, then $\sum_a w_a g_a = 0$, so for any $\delta$, $\sum_a w_a (g_a^\top \delta) = 0$: the first-order effect of any perturbation, pooled across groups at their population weights, must cancel. With exactly two groups this forces $g_a^\top\delta$ and $g_b^\top\delta$ to carry opposite signs; with more than two groups it constrains only the weighted sum, so some groups' first-order terms can share a sign as long as others offset them. Either way, whether $\delta$'s first-order effect on group $a$ is positive or negative is determined by $\delta$'s direction relative to $g_a$, which differs across mechanisms with different $\delta$. This is an empirical question Section~\ref{sec:results} tests directly, not a further consequence of stationarity alone. (Weight decay, or evaluating $g_a$ on held-out data rather than the training objective whose stationary point $\theorig$ actually is, each break the hypothesis of this proposition; our training setup avoids the former, and we treat the latter as using Proposition~\ref{prop:stationarity} as a motivating heuristic for an empirical predictor, not as a literal instantiation of its premise.)
\end{proposition}
This complements Observation~\ref{obs:mechanism-agnostic}: stationarity gives a structural reason to expect compression to reallocate loss between groups to first order whenever $\delta$ is not orthogonal to every $g_a$, while the open, mechanism-dependent question is which group receives that reallocation under a given $\delta$.

\subsection{A Gauss--Newton view: when gradient and curvature decouple}
\label{sec:gauss-newton}

The gradient $g_i$ of a cross-entropy loss for example $i$ is exact, not an approximation; it is the \emph{curvature} that admits a tractable approximation. Writing $p_i$ for the model's predicted probability of example $i$'s true class $y_i$, and $J_i$ for the Jacobian of the pre-softmax logit (for the true class) with respect to $\theta$, the exact per-example gradient satisfies $\|g_i\| = |p_i - 1|\cdot\|J_i\|$, and its \emph{Gauss--Newton} curvature approximation, which discards the term in the exact Hessian weighted by the residual $(p_i-1)$ and the logit's own second derivative, is $p_i(1-p_i)\cdot\|J_i\|^2$. We recover, as a one-line identity, the gradient-curvature relationship \citet{tran2022pruning} establish analytically via distance to the decision boundary, under the Gauss--Newton approximation specifically. This identity also predicts where the approximation's gradient and curvature terms must \emph{decouple}: as $p_i \to 0$ (a confidently misclassified example), $|p_i - 1|$ stays large while $p_i(1-p_i)\to 0$ vanishes. A group dominated by confidently misclassified examples sits in this decoupling regime for the \emph{Gauss--Newton} curvature specifically. The exact Hessian we use empirically (Section~\ref{sec:setup}) retains the residual term the Gauss--Newton approximation drops, and that term is weighted by $(p_i-1)$, which is largest in exactly this regime. Whether the exact-Hessian predictor also degrades there is an empirical question we do not resolve analytically; we treat this subsection as motivation for where the approximation underlying Section~\ref{sec:isotropic-special-case}'s isotropic baseline may be least trustworthy, not as a proof about the full directional predictor. For multi-class losses the Gauss--Newton curvature is $J_i^\top(\mathrm{diag}(p) - pp^\top)J_i$, and the same decoupling occurs as $p$ approaches any vertex of the simplex other than $e_{y_i}$.

\subsection{Isotropic perturbations recover prior quantization theory}
\label{sec:isotropic-special-case}

If a compression mechanism's perturbation is modeled as isotropic, zero-mean noise (as is approximately true of post-training quantization, which rounds each weight independently), then to first order $\mathbb{E}[g_a^\top \delta] = 0$, and to second order $\mathbb{E}[\delta^\top H_a \delta] \approx \sigma^2 \,\mathrm{tr}(H_a)$ for some noise scale $\sigma^2$. The resulting predictor, $\tfrac12\sigma^2\,\mathrm{tr}(H_a)$, is the quantity underlying the trace-based accounts used by \citet{liu2025understanding} and \citet{bellam2024causes,bellam2025explaining} for quantization-specific disparate impact. This predictor and the norm bound of Section~\ref{sec:background} share a structural property our directional predictor does not: both are functions of $\delta$ only through the scalar $\|\delta\|$ (via $\sigma^2 = \|\delta\|^2/d$), so at matched $\|\delta\|$ both are, like the norm bound, insensitive to which mechanism produced $\delta$. The directional theory of Section~\ref{sec:theory} does not compete with the isotropic account; it generalizes it, reducing to it when $\delta$ carries no exploitable directional structure, and the comparison we are actually testing in Section~\ref{sec:results} is between predictors that see only $\|\delta\|$ (norm bound, isotropic) and predictors that see $\delta$'s shape (diagonal, full directional).

\section{Experimental Setup}
\label{sec:setup}

\textbf{Benchmark.} We require a setting with a genuine, precisely controllable spurious correlation and a small, identifiable minority group, in the style of Waterbirds \citep{sagawa2020group} or CelebA \citep{liu2015celeba}. To set the correlation strength and group sizes exactly, rather than inherit whatever a fixed existing dataset provides, we construct an analogous benchmark on CIFAR-10 \citep{krizhevsky2009cifar}: a binary cat-vs-dog task with a painted border color (red/blue) as the spurious feature, correlated with class label for $85\%$ of each class (the majority groups) and \emph{anti}-correlated for the remaining $15\%$ (the minority groups): four groups in total (class $\times$ border color), the same $2\times2$ class-by-spurious-attribute structure as Waterbirds. At this correlation strength the dense model's minority-group accuracy is still low (next paragraph); we found no correlation strength at which minority accuracy is substantially higher while the spurious feature remains strong enough to be learned, so the benchmark does not avoid low-confidence examples, and we do not claim otherwise.

\textbf{Model and training.} We train a CIFAR-style ResNet-18 \citep{he2016resnet} (Section~\ref{sec:background}'s $\theorig$), without weight decay (Proposition~\ref{prop:stationarity} requires a stationary point of the \emph{unregularized} loss), for 12 epochs with SGD, across 3 random seeds.

\textbf{Compression mechanisms.} We implement four mechanisms, matching Section~\ref{sec:theory}: \emph{unstructured} magnitude pruning (global smallest-magnitude threshold across all layers); \emph{structured} channel pruning (smallest-$\ell_2$-norm output channels, ranked \emph{globally} across all layers rather than independently per layer); \emph{low-rank} truncation (SVD of each reshaped conv/linear weight, discarding the globally smallest singular values across all layers); and \emph{quantization} (uniform scalar rounding of every compressible weight to the nearest multiple of a single global step size, Appendix~\ref{app:quantization}). Ranking channels and singular values globally, rather than applying the same severity fraction independently to every layer, matters: with a fixed per-layer fraction, any layer with fewer channels (or lower rank) than roughly the reciprocal of that fraction has its prune count round down to zero, so at low severity compression collapses onto whichever single layer has the most channels or highest rank, leaving the rest of the network untouched. This is not a meaningful stand-in for how either mechanism is used in practice. Quantization's severity parameter is instead a step size in weight units, not a fraction, since that is what the calibration procedure below needs to sweep. No fine-tuning follows compression, so that the measured excess loss reflects the compression perturbation itself rather than subsequent re-optimization.

\textbf{Calibration.} We calibrate each mechanism's severity so that all three produce the same relative weight-perturbation norm, $\|\delta\|/\|\theorig\| \approx 3\%$, via a grid search with local bisection refinement. We emphasize this choice because an earlier attempt to calibrate mechanisms by matching \emph{overall population test loss} instead, a seemingly equivalent and arguably more standard choice, silently permitted much larger severities for some mechanisms, since majority-group accuracy is robust enough that even a large perturbation barely moves population loss. Matching on $\|\delta\|$ directly is the calibration that Observation~\ref{obs:mechanism-agnostic}'s test actually requires (all mechanisms perturbing the model by the same amount), and keeping that perturbation small is what the second-order approximation in Section~\ref{sec:theory} requires to be accurate at all.

\textbf{Predictors compared.} For each (seed, mechanism, group) triple we compute five predictors of $a$'s excess loss under that mechanism's $\delta$: (i) \textbf{full directional}, $g_a^\top\delta + \tfrac12\delta^\top H_a\delta$, via Hessian-vector products \citep{pearlmutter1994fast} (forward-over-reverse-mode automatic differentiation, no Hessian materialization); (ii) \textbf{diagonal}, $g_a^\top\delta + \tfrac12\sum_i H_{a,ii}\delta_i^2$, with $\mathrm{diag}(H_a)$ estimated via Hutchinson's estimator \citep{hutchinson1989trace} (100 Rademacher probes, fixed seed per group for reproducibility), a gradient-retaining variant of the diagonal saliency scores in Section~\ref{sec:related} (real prior methods drop the gradient term entirely); (iii) \textbf{first-order only}, $g_a^\top\delta$, included to isolate whether any advantage of (i) or (ii) comes from curvature at all or only from the correctly signed first-order direction; (iv) \textbf{isotropic}, $\|g_a\|\|\delta\| + \tfrac12\,\mathrm{tr}(H_a)\|\delta\|^2/d$ (Section~\ref{sec:isotropic-special-case}, $d$ the parameter count, trace estimated the same way as (ii); the first term is the same sign-free Cauchy--Schwarz bound used in the norm bound (v), since the isotropic account of \citet{liu2025understanding,bellam2024causes,bellam2025explaining} replaces only the curvature term with a trace, not the first-order term, and that first term, not the trace, turns out to dominate the predictor's behavior empirically); and (v) the \textbf{norm bound}, Eq.~\eqref{eq:norm-bound} itself, $\|g_a\|\|\delta\| + \tfrac12\max(\lambda_{\max}(H_a),0)\|\delta\|^2$, with $\lambda_{\max}$ estimated via power iteration over Hessian-vector products. By Observation~\ref{obs:mechanism-agnostic}, (iv) and (v) depend on $\delta$ only through the scalar $\|\delta\|$ and so predict the same most-harmed group under every mechanism in a given seed; (i)--(iii) depend on $\delta$'s actual coordinates and can vary. All five predictors are computed on a held-out batch of $50$ test examples per group, the same sample size for every group (majority groups have more test examples available but are subsampled to match, so no group's estimate is systematically lower-variance than another's). The \emph{actual} loss change each predictor is scored against is measured separately, as each group's full test-set average loss before and after compression ($850$ examples for the two majority groups, $150$ for the two minority groups), not only the $50$-example predictor batch; the resulting overlap between the predictor's batch and the examples used to measure actual loss is $5.9\%$ for majority groups and $33.3\%$ for minority groups. This overlap is a real, if partial, source of shared sampling noise between predictor and target, distinct from the overfitting risk a learned predictor would carry, since none of the five predictors are fit to data: each is a closed-form function of the gradient and Hessian at $\theorig$, evaluated on a batch, not a score tuned to match actual loss on it. We discuss this further in Section~\ref{sec:discussion}. We also reran the diagonal and isotropic estimators at $4\times$ the probe budget (400 Rademacher probes) as a convergence check, on the original three-mechanism comparison (before quantization was added as a fourth mechanism, Appendix~\ref{app:quantization}): the pooled Spearman correlation for the diagonal predictor was unchanged at the sign-agreement level and, if anything, moved slightly closer to the full-directional and first-order predictors ($\rho=0.888$ vs.\ $\rho=0.876$ at 100 probes), while flip-case accuracy was identical at both probe counts; the isotropic predictor's near-chance performance was unchanged as well, since it is dominated by the norm term rather than the trace term. The comparison in Section~\ref{sec:baseline-comparison} is therefore not an artifact of the probe budget.

\section{Results}
\label{sec:results}

\subsection{Compression mechanisms disagree on who they harm, in every trial}

Across all three random seeds, the four compression mechanisms disagreed on which group suffered the largest increase in test loss: a \emph{flip} in the most-harmed group across mechanisms occurred in $3/3$ seeds (Table~\ref{tab:flips}). Under the norm-based theory (Observation~\ref{obs:mechanism-agnostic}) this should never happen for mechanisms matched on $\|\delta\|$; empirically, it happened every time.

\begin{table}[t]
\centering
\caption{Most-harmed group (by actual test-loss increase) per mechanism and seed. A flip (disagreement across mechanisms) occurs in all three seeds.}
\label{tab:flips}
\begin{tabular}{@{}lcccc@{}}
\toprule
Seed & Unstructured & Channel & Low-rank (SVD) & Quantization \\
\midrule
0 & dog+red (minority)  & dog+red (minority)  & cat+blue (minority) & cat+blue (minority) \\
1 & cat+blue (minority) & dog+blue (majority) & cat+blue (minority) & cat+blue (minority) \\
2 & cat+blue (minority) & dog+red (minority)  & dog+blue (majority) & cat+blue (minority) \\
\bottomrule
\end{tabular}
\end{table}

\subsection{Direction-aware predictors anticipate these flips; magnitude-only predictors do not, but curvature itself adds nothing beyond direction}
\label{sec:baseline-comparison}

\begin{figure}[t]
\centering
% figures/baseline_comparison.tex
% Baseline comparison: which predictor identifies the group a compression
% mechanism harms most? Left: flip-case accuracy (12 flip cases, 4 mechanisms).
% Right: pooled Spearman correlation, predicted vs. actual per-group loss change
% (48 seed x mechanism x group points).
% Source: code/baseline_compare_results.json
%
% Requires in the preamble of main.tex:
%   \usepackage{pgfplots}
%   \pgfplotsset{compat=1.18}        % any compat >= 1.11 works
%   \usetikzlibrary{patterns}
%
% Usage (inside your figure environment):
%   \centering
%   % figures/baseline_comparison.tex
% Baseline comparison: which predictor identifies the group a compression
% mechanism harms most? Left: flip-case accuracy (12 flip cases, 4 mechanisms).
% Right: pooled Spearman correlation, predicted vs. actual per-group loss change
% (48 seed x mechanism x group points).
% Source: code/baseline_compare_results.json
%
% Requires in the preamble of main.tex:
%   \usepackage{pgfplots}
%   \pgfplotsset{compat=1.18}        % any compat >= 1.11 works
%   \usetikzlibrary{patterns}
%
% Usage (inside your figure environment):
%   \centering
%   % figures/baseline_comparison.tex
% Baseline comparison: which predictor identifies the group a compression
% mechanism harms most? Left: flip-case accuracy (12 flip cases, 4 mechanisms).
% Right: pooled Spearman correlation, predicted vs. actual per-group loss change
% (48 seed x mechanism x group points).
% Source: code/baseline_compare_results.json
%
% Requires in the preamble of main.tex:
%   \usepackage{pgfplots}
%   \pgfplotsset{compat=1.18}        % any compat >= 1.11 works
%   \usetikzlibrary{patterns}
%
% Usage (inside your figure environment):
%   \centering
%   \input{figures/baseline_comparison.tex}
%
% Encoding is redundant so the figure works in grayscale/B&W print:
%   full       = dark blue   (Okabe-Ito #0072B2), solid fill
%   diagonal   = sky blue    (Okabe-Ito #56B4E9), diagonal hatching
%   first-order= bluish green(Okabe-Ito #009E73), horizontal hatching
%   isotropic  = orange      (Okabe-Ito #E69F00), crosshatching
%   norm-bound = light gray, dots
% Each bar is also identified by its x-axis label. The three direction-aware
% predictors (full, diagonal, first-order) are statistically indistinguishable
% from one another and are visually grouped together; the two magnitude-only
% predictors (isotropic, norm-bound) are grouped on the right.
\definecolor{bcOurs}{HTML}{0072B2}%
\definecolor{bcDiag}{HTML}{56B4E9}%
\definecolor{bcFirst}{HTML}{009E73}%
\definecolor{bcIso}{HTML}{E69F00}%
\definecolor{bcNorm}{HTML}{BBBBBB}%
\begin{tikzpicture}
  \pgfplotsset{
    bcpanel/.style={
      width=2.85in,
      height=2.05in,
      scale only axis=false,
      ybar,
      bar width=0.27in,
      symbolic x coords={full,diag,first,iso,norm},
      xmin=full, xmax=norm,
      xtick={full,diag,first,iso,norm},
      xticklabels={{Full\\(HVP)},{Diagonal\\Hessian},{First-\\order only},{Isotropic\\(Liu et al.)},{Norm-\\bound}},
      xticklabel style={align=center, font=\scriptsize},
      x tick label style={text height=1.5ex},
      enlarge x limits=0.13,
      tick label style={font=\scriptsize},
      ylabel style={font=\small},
      title style={font=\small, yshift=-1ex},
      ymajorgrids,
      grid style={gray!25, line width=0.3pt},
      major x tick style={draw=none},
      axis line style={gray!70},
      nodes near coords,
      nodes near coords align=vertical,
      point meta=explicit symbolic,
      nodes near coords style={font=\scriptsize, align=center},
      every axis plot/.append style={bar shift=0pt, draw=black, line width=0.5pt},
    },
  }

  % ---------------- Left panel: flip-case accuracy ----------------
  \begin{axis}[
      bcpanel,
      name=acc,
      title={(a) Flip-case accuracy},
      ylabel={Accuracy (\%)},
      ymin=0, ymax=114,
      ytick={0,25,50,75,100},
    ]
    \addplot[fill=bcOurs]
      coordinates {(full,83.3) [{83.3\%\\[-0.3ex](10/12)}]};
    \addplot[fill=bcDiag, postaction={pattern=north east lines}]
      coordinates {(diag,83.3) [{83.3\%\\[-0.3ex](10/12)}]};
    \addplot[fill=bcFirst, postaction={pattern=horizontal lines}]
      coordinates {(first,83.3) [{83.3\%\\[-0.3ex](10/12)}]};
    \addplot[fill=bcIso, postaction={pattern=crosshatch}]
      coordinates {(iso,25.0)  [{25.0\%\\[-0.3ex](3/12)}]};
    \addplot[fill=bcNorm, postaction={pattern=dots}]
      coordinates {(norm,25.0) [{25.0\%\\[-0.3ex](3/12)}]};
    % Chance level: 4 groups -> 25%
    \draw[densely dashed, black!70, line width=0.6pt]
      ({rel axis cs:0,0} |- {axis cs:full,25}) -- ({rel axis cs:1,0} |- {axis cs:full,25});
    % Key for the dashed line, placed in the empty upper-right region
    \draw[densely dashed, black!70, line width=0.6pt]
      ([xshift=-0.30in]axis cs:iso,95) -- ([xshift=-0.05in]axis cs:iso,95)
      node[right, font=\scriptsize, text=black!80] {chance (25\%)};
    % Brace calling out that the three direction-aware predictors tie
    \draw[decorate, decoration={brace, amplitude=4pt, mirror}, black!60]
      (axis cs:full,-6) -- (axis cs:first,-6)
      node[midway, below, yshift=-5pt, font=\scriptsize, align=center, text=black!70]
      {direction-aware: statistically tied};
  \end{axis}

  % ---------------- Right panel: pooled Spearman correlation ----------------
  \begin{axis}[
      bcpanel,
      at={(acc.outer east)}, anchor=outer west, xshift=0.05in,
      title={(b) Spearman $\rho$ (predicted vs.\ actual)},
      ylabel={Spearman $\rho$},
      ymin=-0.4, ymax=0.98,
      ytick={-0.25,0,0.25,0.5,0.75},
      yticklabel style={/pgf/number format/fixed, /pgf/number format/precision=2},
    ]
    \addplot[fill=bcOurs]
      coordinates {(full,0.901) [0.901]};
    \addplot[fill=bcDiag, postaction={pattern=north east lines}]
      coordinates {(diag,0.901) [0.901]};
    \addplot[fill=bcFirst, postaction={pattern=horizontal lines}]
      coordinates {(first,0.904) [0.904]};
    % Negative, non-significant bar: put its value label below the bar end
    \addplot[fill=bcIso, postaction={pattern=crosshatch},
              nodes near coords style={anchor=north}]
      coordinates {(iso,-0.041) [{$-0.041$\\[-0.3ex](n.s.)}]};
    % Zero reference line
    \draw[black!80, line width=0.5pt]
      ({rel axis cs:0,0} |- {axis cs:full,0}) -- ({rel axis cs:1,0} |- {axis cs:full,0});
    % Norm-bound outputs only a group ranking, not per-group loss changes
    \node[font=\scriptsize, align=center, text=black!70, fill=white, inner sep=1.5pt,
          anchor=south] at (axis cs:norm,0.04)
      {n/a\\[-0.3ex](ranking\\[-0.3ex]only)};
  \end{axis}
\end{tikzpicture}%

%
% Encoding is redundant so the figure works in grayscale/B&W print:
%   full       = dark blue   (Okabe-Ito #0072B2), solid fill
%   diagonal   = sky blue    (Okabe-Ito #56B4E9), diagonal hatching
%   first-order= bluish green(Okabe-Ito #009E73), horizontal hatching
%   isotropic  = orange      (Okabe-Ito #E69F00), crosshatching
%   norm-bound = light gray, dots
% Each bar is also identified by its x-axis label. The three direction-aware
% predictors (full, diagonal, first-order) are statistically indistinguishable
% from one another and are visually grouped together; the two magnitude-only
% predictors (isotropic, norm-bound) are grouped on the right.
\definecolor{bcOurs}{HTML}{0072B2}%
\definecolor{bcDiag}{HTML}{56B4E9}%
\definecolor{bcFirst}{HTML}{009E73}%
\definecolor{bcIso}{HTML}{E69F00}%
\definecolor{bcNorm}{HTML}{BBBBBB}%
\begin{tikzpicture}
  \pgfplotsset{
    bcpanel/.style={
      width=2.85in,
      height=2.05in,
      scale only axis=false,
      ybar,
      bar width=0.27in,
      symbolic x coords={full,diag,first,iso,norm},
      xmin=full, xmax=norm,
      xtick={full,diag,first,iso,norm},
      xticklabels={{Full\\(HVP)},{Diagonal\\Hessian},{First-\\order only},{Isotropic\\(Liu et al.)},{Norm-\\bound}},
      xticklabel style={align=center, font=\scriptsize},
      x tick label style={text height=1.5ex},
      enlarge x limits=0.13,
      tick label style={font=\scriptsize},
      ylabel style={font=\small},
      title style={font=\small, yshift=-1ex},
      ymajorgrids,
      grid style={gray!25, line width=0.3pt},
      major x tick style={draw=none},
      axis line style={gray!70},
      nodes near coords,
      nodes near coords align=vertical,
      point meta=explicit symbolic,
      nodes near coords style={font=\scriptsize, align=center},
      every axis plot/.append style={bar shift=0pt, draw=black, line width=0.5pt},
    },
  }

  % ---------------- Left panel: flip-case accuracy ----------------
  \begin{axis}[
      bcpanel,
      name=acc,
      title={(a) Flip-case accuracy},
      ylabel={Accuracy (\%)},
      ymin=0, ymax=114,
      ytick={0,25,50,75,100},
    ]
    \addplot[fill=bcOurs]
      coordinates {(full,83.3) [{83.3\%\\[-0.3ex](10/12)}]};
    \addplot[fill=bcDiag, postaction={pattern=north east lines}]
      coordinates {(diag,83.3) [{83.3\%\\[-0.3ex](10/12)}]};
    \addplot[fill=bcFirst, postaction={pattern=horizontal lines}]
      coordinates {(first,83.3) [{83.3\%\\[-0.3ex](10/12)}]};
    \addplot[fill=bcIso, postaction={pattern=crosshatch}]
      coordinates {(iso,25.0)  [{25.0\%\\[-0.3ex](3/12)}]};
    \addplot[fill=bcNorm, postaction={pattern=dots}]
      coordinates {(norm,25.0) [{25.0\%\\[-0.3ex](3/12)}]};
    % Chance level: 4 groups -> 25%
    \draw[densely dashed, black!70, line width=0.6pt]
      ({rel axis cs:0,0} |- {axis cs:full,25}) -- ({rel axis cs:1,0} |- {axis cs:full,25});
    % Key for the dashed line, placed in the empty upper-right region
    \draw[densely dashed, black!70, line width=0.6pt]
      ([xshift=-0.30in]axis cs:iso,95) -- ([xshift=-0.05in]axis cs:iso,95)
      node[right, font=\scriptsize, text=black!80] {chance (25\%)};
    % Brace calling out that the three direction-aware predictors tie
    \draw[decorate, decoration={brace, amplitude=4pt, mirror}, black!60]
      (axis cs:full,-6) -- (axis cs:first,-6)
      node[midway, below, yshift=-5pt, font=\scriptsize, align=center, text=black!70]
      {direction-aware: statistically tied};
  \end{axis}

  % ---------------- Right panel: pooled Spearman correlation ----------------
  \begin{axis}[
      bcpanel,
      at={(acc.outer east)}, anchor=outer west, xshift=0.05in,
      title={(b) Spearman $\rho$ (predicted vs.\ actual)},
      ylabel={Spearman $\rho$},
      ymin=-0.4, ymax=0.98,
      ytick={-0.25,0,0.25,0.5,0.75},
      yticklabel style={/pgf/number format/fixed, /pgf/number format/precision=2},
    ]
    \addplot[fill=bcOurs]
      coordinates {(full,0.901) [0.901]};
    \addplot[fill=bcDiag, postaction={pattern=north east lines}]
      coordinates {(diag,0.901) [0.901]};
    \addplot[fill=bcFirst, postaction={pattern=horizontal lines}]
      coordinates {(first,0.904) [0.904]};
    % Negative, non-significant bar: put its value label below the bar end
    \addplot[fill=bcIso, postaction={pattern=crosshatch},
              nodes near coords style={anchor=north}]
      coordinates {(iso,-0.041) [{$-0.041$\\[-0.3ex](n.s.)}]};
    % Zero reference line
    \draw[black!80, line width=0.5pt]
      ({rel axis cs:0,0} |- {axis cs:full,0}) -- ({rel axis cs:1,0} |- {axis cs:full,0});
    % Norm-bound outputs only a group ranking, not per-group loss changes
    \node[font=\scriptsize, align=center, text=black!70, fill=white, inner sep=1.5pt,
          anchor=south] at (axis cs:norm,0.04)
      {n/a\\[-0.3ex](ranking\\[-0.3ex]only)};
  \end{axis}
\end{tikzpicture}%

%
% Encoding is redundant so the figure works in grayscale/B&W print:
%   full       = dark blue   (Okabe-Ito #0072B2), solid fill
%   diagonal   = sky blue    (Okabe-Ito #56B4E9), diagonal hatching
%   first-order= bluish green(Okabe-Ito #009E73), horizontal hatching
%   isotropic  = orange      (Okabe-Ito #E69F00), crosshatching
%   norm-bound = light gray, dots
% Each bar is also identified by its x-axis label. The three direction-aware
% predictors (full, diagonal, first-order) are statistically indistinguishable
% from one another and are visually grouped together; the two magnitude-only
% predictors (isotropic, norm-bound) are grouped on the right.
\definecolor{bcOurs}{HTML}{0072B2}%
\definecolor{bcDiag}{HTML}{56B4E9}%
\definecolor{bcFirst}{HTML}{009E73}%
\definecolor{bcIso}{HTML}{E69F00}%
\definecolor{bcNorm}{HTML}{BBBBBB}%
\begin{tikzpicture}
  \pgfplotsset{
    bcpanel/.style={
      width=2.85in,
      height=2.05in,
      scale only axis=false,
      ybar,
      bar width=0.27in,
      symbolic x coords={full,diag,first,iso,norm},
      xmin=full, xmax=norm,
      xtick={full,diag,first,iso,norm},
      xticklabels={{Full\\(HVP)},{Diagonal\\Hessian},{First-\\order only},{Isotropic\\(Liu et al.)},{Norm-\\bound}},
      xticklabel style={align=center, font=\scriptsize},
      x tick label style={text height=1.5ex},
      enlarge x limits=0.13,
      tick label style={font=\scriptsize},
      ylabel style={font=\small},
      title style={font=\small, yshift=-1ex},
      ymajorgrids,
      grid style={gray!25, line width=0.3pt},
      major x tick style={draw=none},
      axis line style={gray!70},
      nodes near coords,
      nodes near coords align=vertical,
      point meta=explicit symbolic,
      nodes near coords style={font=\scriptsize, align=center},
      every axis plot/.append style={bar shift=0pt, draw=black, line width=0.5pt},
    },
  }

  % ---------------- Left panel: flip-case accuracy ----------------
  \begin{axis}[
      bcpanel,
      name=acc,
      title={(a) Flip-case accuracy},
      ylabel={Accuracy (\%)},
      ymin=0, ymax=114,
      ytick={0,25,50,75,100},
    ]
    \addplot[fill=bcOurs]
      coordinates {(full,83.3) [{83.3\%\\[-0.3ex](10/12)}]};
    \addplot[fill=bcDiag, postaction={pattern=north east lines}]
      coordinates {(diag,83.3) [{83.3\%\\[-0.3ex](10/12)}]};
    \addplot[fill=bcFirst, postaction={pattern=horizontal lines}]
      coordinates {(first,83.3) [{83.3\%\\[-0.3ex](10/12)}]};
    \addplot[fill=bcIso, postaction={pattern=crosshatch}]
      coordinates {(iso,25.0)  [{25.0\%\\[-0.3ex](3/12)}]};
    \addplot[fill=bcNorm, postaction={pattern=dots}]
      coordinates {(norm,25.0) [{25.0\%\\[-0.3ex](3/12)}]};
    % Chance level: 4 groups -> 25%
    \draw[densely dashed, black!70, line width=0.6pt]
      ({rel axis cs:0,0} |- {axis cs:full,25}) -- ({rel axis cs:1,0} |- {axis cs:full,25});
    % Key for the dashed line, placed in the empty upper-right region
    \draw[densely dashed, black!70, line width=0.6pt]
      ([xshift=-0.30in]axis cs:iso,95) -- ([xshift=-0.05in]axis cs:iso,95)
      node[right, font=\scriptsize, text=black!80] {chance (25\%)};
    % Brace calling out that the three direction-aware predictors tie
    \draw[decorate, decoration={brace, amplitude=4pt, mirror}, black!60]
      (axis cs:full,-6) -- (axis cs:first,-6)
      node[midway, below, yshift=-5pt, font=\scriptsize, align=center, text=black!70]
      {direction-aware: statistically tied};
  \end{axis}

  % ---------------- Right panel: pooled Spearman correlation ----------------
  \begin{axis}[
      bcpanel,
      at={(acc.outer east)}, anchor=outer west, xshift=0.05in,
      title={(b) Spearman $\rho$ (predicted vs.\ actual)},
      ylabel={Spearman $\rho$},
      ymin=-0.4, ymax=0.98,
      ytick={-0.25,0,0.25,0.5,0.75},
      yticklabel style={/pgf/number format/fixed, /pgf/number format/precision=2},
    ]
    \addplot[fill=bcOurs]
      coordinates {(full,0.901) [0.901]};
    \addplot[fill=bcDiag, postaction={pattern=north east lines}]
      coordinates {(diag,0.901) [0.901]};
    \addplot[fill=bcFirst, postaction={pattern=horizontal lines}]
      coordinates {(first,0.904) [0.904]};
    % Negative, non-significant bar: put its value label below the bar end
    \addplot[fill=bcIso, postaction={pattern=crosshatch},
              nodes near coords style={anchor=north}]
      coordinates {(iso,-0.041) [{$-0.041$\\[-0.3ex](n.s.)}]};
    % Zero reference line
    \draw[black!80, line width=0.5pt]
      ({rel axis cs:0,0} |- {axis cs:full,0}) -- ({rel axis cs:1,0} |- {axis cs:full,0});
    % Norm-bound outputs only a group ranking, not per-group loss changes
    \node[font=\scriptsize, align=center, text=black!70, fill=white, inner sep=1.5pt,
          anchor=south] at (axis cs:norm,0.04)
      {n/a\\[-0.3ex](ranking\\[-0.3ex]only)};
  \end{axis}
\end{tikzpicture}%

\caption{(a) Accuracy at identifying the actual most-harmed group, restricted to the $n=12$ \emph{flip cases}: (seed, mechanism) pairs from Table~\ref{tab:flips} where mechanisms disagree on which group suffers most, the hardest and most diagnostic test of a predictor. Dashed line: chance accuracy for a 4-group benchmark ($25\%$). Norm-bound has no $\rho$ in (b) because, by Observation~\ref{obs:mechanism-agnostic}, it predicts the same single group regardless of mechanism rather than a per-group loss value, so it cannot be correlated against the actual per-group loss change. (b) Pooled Spearman correlation between each predictor's score and the actual per-group test-loss change, across all $n=48$ (seed $\times$ mechanism $\times$ group) observations. The three \emph{direction-aware} predictors (full directional, Section~\ref{sec:theory}; diagonal-only, used in prior mitigation methods; and the bare first-order term alone) are statistically indistinguishable from one another and all substantially outperform the two \emph{magnitude-only} predictors (isotropic/trace-based, Section~\ref{sec:isotropic-special-case}; and the norm bound), whose correlation is statistically indistinguishable from zero ($p=0.78$) and whose flip-case accuracy matches chance exactly.}
\label{fig:baseline-comparison}
\end{figure}

Figure~\ref{fig:baseline-comparison} compares all five predictors from Section~\ref{sec:setup}, now across all four mechanisms including quantization. All three \emph{direction-aware} predictors identify the actual most-harmed group in $10/12$ ($83.3\%$) of the flip cases from Table~\ref{tab:flips}: the full directional predictor (pooled Spearman $\rho=0.901$, $p=2.4\times10^{-18}$, $n=48$), the diagonal-only predictor ($\rho=0.901$, $p=2.9\times10^{-18}$), and, the more surprising result, the bare first-order term $g_a^\top\delta$ alone, with \emph{no} curvature information at all ($\rho=0.904$, $p=1.4\times10^{-18}$). These three accuracies and three correlations are statistically indistinguishable from one another. On this benchmark, curvature, whether the full off-diagonal Hessian or only its diagonal, adds no measurable predictive power beyond simply getting the sign and magnitude of the first-order term right. This qualifies our motivation in Sections~\ref{sec:related} and~\ref{sec:theory} for comparing curvature-aware predictors in the first place: the comparison we set out to run (full vs.\ diagonal vs.\ isotropic) turns out to be dominated by a plainer one, direction versus no direction at all, which is exactly what Proposition~\ref{prop:stationarity}'s purely first-order argument already predicts should matter. We also find this echoes what the classical, fairness-agnostic pruning-saliency literature has found for pooled loss \citep{lee2019snip,molchanov2019importance}: a first-order criterion alone is often enough, and curvature earns its cost less often than its prominence in that literature, and in the group-conditioned mitigation methods of Section~\ref{sec:related}, would suggest. The decisive contrast remains between \emph{any} direction-aware predictor and the two predictors built only from $\|\delta\|$: the isotropic predictor, at $3/12$ ($25.0\%$, exactly chance) with a Spearman correlation of $-0.041$ that is \emph{not} statistically distinguishable from zero ($p=0.78$, $n=48$), and the mechanism-independent norm bound, also at $3/12$ ($25.0\%$) by construction (Observation~\ref{obs:mechanism-agnostic}).

Quantization's predicted-vs-actual values from the full directional predictor track especially tightly (e.g.\ seed 0, group cat+blue: actual $+0.0143$, predicted $+0.0146$; Appendix~\ref{app:quantization} gives the full per-seed numbers). This is consistent with Section~\ref{sec:isotropic-special-case}'s claim that quantization's perturbation is one the second-order expansion describes well. We also tested that section's more specific prediction directly: because quantization's perturbation is close to isotropic noise, the trace-based curvature approximation $\mathrm{tr}(H_a)\|\delta\|^2/d$ should match the true quadratic term $\delta^\top H_a\delta$ more closely for quantization than for the other three mechanisms. We hold the comparison to the curvature term alone, since the full isotropic predictor is dominated by its sign-blind first-order term (Appendix~\ref{app:quantization}). The result only partly confirms this: quantization has the lowest typical relative error between the two quantities (median $65\%$, versus $95$--$157\%$ for the other three mechanisms), but at our sample size ($n=12$ per mechanism, 3 seeds $\times$ 4 groups) the correlation between true and approximated curvature does not cleanly single out quantization as the best-matched mechanism; low-rank truncation's correlation is in fact higher. Appendix~\ref{app:quantization} gives the full numbers and what a larger sample would need to resolve. The isotropic score's \emph{overall} flip-case accuracy sits at chance regardless, no better on quantization than on the other three mechanisms: its sign-blind first-order term dominates its behavior by construction (Section~\ref{sec:setup}), independent of whether any one mechanism's perturbation happens to be well-approximated as isotropic. We take the direction-vs-magnitude result as the paper's central empirical finding, and the near-equality among direction-aware predictors as a genuine, more modest finding alongside it. Section~\ref{sec:discussion} discusses its practical upshot and its scope as a limitation. Throughout, we report loss rather than accuracy, since loss is the continuous quantity the second-order theory predicts directly; Appendix~\ref{app:accuracy-effect-sizes} reports the corresponding accuracy effect sizes at each seed's most-harmed group and finds they are often small in absolute terms, not because the loss predictions are wrong but because many minority groups already sit near floor accuracy before any compression is applied (Section~\ref{sec:setup}).

\subsection{A calibration pitfall, and an open question about the valid regime}
\label{sec:validity-frontier}

Our first attempt at this comparison calibrated mechanisms by matching \emph{overall population test loss} rather than $\|\delta\|$ directly (Section~\ref{sec:setup}). Because majority-group loss is robust to large perturbations, this silently selected much larger severities for some mechanisms than for others while still appearing matched on the only quantity it controlled (population loss), and the resulting predictions were far less reliable, including outright sign errors. We do not report the specific figures from that attempt, since it used an earlier, flawed implementation of the channel-pruning and low-rank mechanisms (Section~\ref{sec:setup}) and is not a clean ablation of perturbation size alone; we report the methodological lesson instead, since it generalizes beyond our specific implementation: \emph{matching mechanisms on their effect on overall loss is not the same as matching the perturbation each one actually makes}, and only the latter is the condition any second-order account of compression requires. We ran exactly this sweep with the corrected mechanism implementations (Appendix~\ref{app:severity-sweep}), across four severity targets from $1\%$ to $25\%$ relative perturbation. The full directional predictor's pooled correlation with actual loss does not monotonically degrade with severity ($\rho=0.997$ at $1\%$, down to $\rho=0.901$ at $3\%$, then back up to $\rho\approx0.97$--$0.98$ at $10\%$ and $25\%$), but its flip-case accuracy, the harder and more practically relevant test, does: from a perfect $100\%$ at $1\%$ severity down to $58.3\%$ at $25\%$, approaching but not reaching chance. The isotropic predictor, meanwhile, becomes progressively more anti-correlated with actual loss as severity grows ($\rho=0.096$ at $1\%$ to $\rho=-0.337$ at $25\%$), rather than merely remaining uninformative. Our Waterbirds replication (Appendix~\ref{app:waterbirds}) speaks to a second axis of the same question that is not $\|\delta\|$ at all: holding severity fixed at $3\%$ and varying only how overfit the pre-compression model is, we find the same qualitative transition, from the gate failing outright on an over-trained checkpoint to it passing on an earlier-stopped one. This suggests the valid regime is governed jointly by perturbation size and by how far a group's loss already sits from the point where the local expansion in Section~\ref{sec:theory} is taken, not by perturbation size alone.

\section{Discussion}
\label{sec:discussion}

\textbf{Mechanism choice is a fairness decision.} We show that a practitioner's choice of \emph{which} compression technique to deploy, not only how aggressively to compress, changes which subpopulation absorbs the resulting loss increase. A direction-aware predictor requires $\delta$, so it cannot replace measuring actual per-group loss on an already-compressed model when that measurement is possible; its practical value is instead in \emph{amortizing} access to group-labeled data, which is often restricted or sensitive to collect. Section~\ref{sec:baseline-comparison}'s finding that curvature adds no measurable benefit here makes this amortization cheaper than we originally set out to show: computing only $g_a$ once per group, under whatever audited or privileged access produced the group labels, lets a practitioner score any number of \emph{candidate} compression mechanisms with a single gradient inner product per (group, candidate), no Hessian-vector product and no $H_a$ at all, without repeating that data access for each one, and before committing to deploy any of them.

\textbf{Limitations.} Our main empirical validation uses a controlled substitute for standard spurious-correlation benchmarks (Section~\ref{sec:setup}), built to let us set the correlation strength and group structure precisely rather than to work around any single constraint. We replicate it on the real benchmark it was built to approximate, Waterbirds (Appendix~\ref{app:waterbirds}), and the direction-vs-magnitude result holds there too under a standard early-stopped training recipe ($77.8\%$ flip accuracy, matching the substitute's own figure almost exactly), but a more aggressively-trained checkpoint on the same data and mechanisms fails the gate entirely, with even the direction-vs-magnitude distinction disappearing. This is informative about the open question below, not only a caveat: the real-dataset replication indicates that a model's degree of overfitting at $\theorig$, not only compression severity, governs the second-order account's validity. Our evidence base is $3$ seeds per dataset; a larger-scale replication would strengthen confidence in the exact accuracy figures. The near-equality we find among the full, diagonal, and first-order-only predictors (Section~\ref{sec:baseline-comparison}) may be specific to the mild severity regime we test ($\|\delta\|/\|\theta^\star\|\approx3\%$): the quadratic term is small relative to the linear one by construction in this regime, so whether curvature's contribution grows at higher severity, where Section~\ref{sec:validity-frontier} already shows the second-order theory itself begins to degrade, is an open question we have not resolved, and the two questions (does curvature help, and is the second-order theory still valid) may need to be answered jointly rather than separately. We have included quantization as a fourth mechanism in the main comparison (Section~\ref{sec:baseline-comparison}); it strengthens the headline gate numbers, but the more specific confirmatory test it enables for Section~\ref{sec:isotropic-special-case}'s isotropic-special-case prediction (that quantization's curvature term specifically should be better-approximated by the trace-based isotropic account than the other three mechanisms') comes out as a partial, sample-size-limited result rather than a clean confirmation (Appendix~\ref{app:quantization}); a larger sample would be needed to resolve it further. We also match mechanisms on parameter-space weight-perturbation norm, which is not identical to matching their effect on the model's input-output function: batch normalization makes preceding convolutional weights scale-invariant during training, so this norm is not perfectly comparable across mechanisms that touch different layers. A function-space calibration (e.g.\ matched change in output logits) is a more faithful but more expensive alternative we leave for future work. Our quantization mechanism (Section~\ref{sec:setup}) is a uniform scalar rounding at a single global step size, chosen to match the other three mechanisms' no-mechanism-specific-tuning treatment; it does not represent the activation-aware or per-channel schemes used in practice for large models (e.g.\ AWQ \citep{lin2024awq}, GPTQ \citep{frantar2023gptq}), which impose their own directional structure on $\delta$ and may interact with group-specific harm differently than uniform rounding does. Our main training recipe uses no weight decay, required by Proposition~\ref{prop:stationarity}'s premise. Appendix~\ref{app:weight-decay} tests whether the empirical findings themselves, as opposed to the proposition's analytic guarantee, survive standard $\ell_2$ regularization, by retraining with weight decay $5\times10^{-4}$: the gate still passes cleanly (pooled Spearman $\rho=0.975$, $3/3$ flip seeds, $75.0\%$ flip-case accuracy), with the same direction-vs-magnitude pattern and the same near-equality among direction-aware predictors. We rank channels and singular values globally across layers rather than enforcing a fixed per-layer budget (Section~\ref{sec:setup} explains why: a fixed per-layer fraction can collapse severity onto a single layer); whether our findings are specific to global ranking or also hold under other practical budgeting schemes is likewise untested. Finally, the predictors in Section~\ref{sec:setup} are evaluated on a $50$-example batch per group that partially overlaps with the larger set used to measure actual loss ($5.9\%$ overlap for majority groups, $33.3\%$ for minority groups). Because every predictor is a closed-form function of $\theorig$, not a score fit to data, this overlap cannot let a predictor memorize test-set idiosyncrasies the way a learned model could, but it can still let predictor and target share sampling noise, especially for the minority groups where the overlap is largest. We reran the comparison with fully disjoint predictor and evaluation samples per group (Appendix~\ref{app:disjoint-split}) and recovered nearly identical results, confirming this overlap is not driving the main finding.

\section{Conclusion}
\label{sec:conclusion}

Which group pays the accuracy cost of model compression is not a fixed property of the data or the model alone. It is a joint property of the model's loss landscape and the specific geometric shape of a given compression mechanism's weight perturbation. Existing disparate-impact theory, built on perturbation magnitude alone, cannot see this and predicts, incorrectly, that mechanism choice should not matter. We show it does, and that anticipating it requires only perturbation \emph{direction}, not curvature: a predictor built from the bare first-order term performs on par with one using the full Hessian, so a single gradient per group, no Hessian-vector product at all, gives a cheap diagnostic that anticipates mechanism-specific harm within a well-characterized regime. We map where that regime ends.

% cpal_2027.sty already sets \bibliographystyle{unsrtnat} when natbib loads
\bibliography{refs}

\appendix
\section{Appendix}

\subsection{Per-seed calibration details}
\label{app:calibration}
Table~\ref{tab:calibration} reports, for each (seed, mechanism) pair used in the main results (Section~\ref{sec:results}), the severity parameter found by the grid-search-plus-bisection calibration procedure (Section~\ref{sec:setup}) and the resulting achieved relative weight-perturbation norm $\|\delta\|/\|\theorig\|$. For the first three mechanisms the severity parameter is the global fraction of components removed (individual weights for unstructured pruning, output channels for channel pruning, singular values for low-rank truncation), ranked across all compressible layers jointly, not per layer (Section~\ref{sec:setup}); for quantization it is instead the global rounding step size in weight units, not a fraction. The four mechanisms require very different severities to reach the same $\|\delta\|/\|\theorig\|\approx3\%$ target: unstructured pruning removes roughly $11$--$12\%$ of individual weights, channel pruning and low-rank truncation reach the same perturbation norm by removing well under $1\%$ of channels or singular values respectively (each removed channel or singular value concentrates much more of the weight matrix's mass than a single pruned weight), and quantization needs a step size roughly three orders of magnitude smaller than the typical weight magnitude, consistent with uniform rounding perturbing every compressible weight at once rather than a sparse subset. All four mechanisms' achieved norms fall within $2.91$--$3.07\%$ of the $3\%$ target across all $12$ (seed, mechanism) pairs. This is the calibration that avoids the population-loss-matching pitfall discussed in Section~\ref{sec:validity-frontier}: matching on $\|\delta\|$ directly, rather than on overall test loss, is what lets severity differ this much across mechanisms while keeping the quantity Observation~\ref{obs:mechanism-agnostic} and Section~\ref{sec:theory} actually require, the perturbation size each mechanism makes to the model, fixed.

\begin{table}[h]
\centering
\caption{Per-seed, per-mechanism calibration. \emph{Severity} is the global fraction of components removed for the first three mechanisms, and a global rounding step size in weight units for quantization (Appendix~\ref{app:calibration}); \emph{achieved $\|\delta\|/\|\theorig\|$} is the resulting relative weight-perturbation norm, targeting $3\%$.}
\label{tab:calibration}
\begin{tabular}{@{}llrr@{}}
\toprule
Seed & Mechanism & Severity & Achieved $\|\delta\|/\|\theorig\|$ \\
\midrule
\multirow{4}{*}{0} & Unstructured & $0.1169$ & $2.999\%$ \\
 & Channel & $0.00150$ & $3.069\%$ \\
 & Low-rank (SVD) & $0.00494$ & $2.981\%$ \\
 & Quantization & $0.00137$ & $3.000\%$ \\
\addlinespace
\multirow{4}{*}{1} & Unstructured & $0.1148$ & $3.000\%$ \\
 & Channel & $0.00132$ & $2.913\%$ \\
 & Low-rank (SVD) & $0.00494$ & $2.969\%$ \\
 & Quantization & $0.00133$ & $2.999\%$ \\
\addlinespace
\multirow{4}{*}{2} & Unstructured & $0.1126$ & $2.999\%$ \\
 & Channel & $0.00132$ & $2.964\%$ \\
 & Low-rank (SVD) & $0.00451$ & $2.977\%$ \\
 & Quantization & $0.00129$ & $2.999\%$ \\
\bottomrule
\end{tabular}
\end{table}

\subsection{Waterbirds replication}
\label{app:waterbirds}

We replicate the comparison in Section~\ref{sec:results} on Waterbirds \citep{sagawa2020group} (\texttt{waterbird\_complete95\_forest2water2}), the real dataset the CIFAR-10 substitute (Section~\ref{sec:setup}) was built to approximate: a bird cropped from CUB-200-2011 \citep{wah2011caltech} and composited onto a land or water background from the Places dataset \citep{zhou2017places}, binary label (landbird/waterbird) spuriously correlated with background at $95\%$ in training ($4795$ images; the smallest group, waterbird-on-land, has only $56$) and group-balanced in testing ($5794$ images). Groups follow the standard convention $\mathrm{group}=2\cdot\mathrm{label}+\mathrm{place}$. The original CodaLab-hosted download for this dataset was unavailable at the time of writing; we used an independently-hosted mirror of the identical, officially-released split instead. We fine-tune an ImageNet-pretrained ResNet-18 end-to-end (images resized to $128^2$, SGD, no weight decay, $3$ seeds), otherwise following Section~\ref{sec:setup}'s procedure unchanged: the same three mechanisms, the same global-ranking calibration to $\|\delta\|/\|\theorig\|\approx3\%$, and the same gate.

\textbf{A standard training protocol fails the gate.} Training to convergence by the usual criterion (loss plateaued, here $8$ epochs, reaching $100\%$ \emph{training} accuracy in every seed) fails the gate: pooled Spearman $\rho=0.559$ ($<0.80$) and, on the $2/3$ seeds that do flip, directional accuracy is only $3/6=50.0\%$ ($<75\%$). More strikingly, \emph{every} predictor, including the magnitude-only isotropic and norm-bound baselines which Section~\ref{sec:baseline-comparison} found at chance on the CIFAR substitute, lands at exactly the same $50.0\%$ on this checkpoint: the direction-vs-magnitude distinction that is the paper's central finding disappears entirely, not just the weaker curvature-vs-no-curvature one. Base test loss on the smallest group (waterbird-on-land) is $\approx\!2.9$ nats, i.e.\ the model is confidently \emph{wrong} on most of that group before any compression is applied, a regime our theory does not claim to cover, since Section~\ref{sec:theory}'s Taylor expansion is a \emph{local} approximation around $\theorig$ and places no guarantee on its accuracy once a group's loss is already far from where the expansion is taken.

\textbf{Earlier stopping restores it.} As a direct test of the hypothesis that this failure is a symptom of overfitting rather than of the dataset itself, we re-ran all $3$ seeds stopped at epoch $3$ (training accuracy $99.3$--$99.5\%$, not yet saturated), the only change from the $8$-epoch run. This single change is enough: pooled Spearman rises to $\rho=0.747$ (still short of the $0.80$ gate), all $3$ seeds now produce flips, and directional accuracy on those $9$ flip cases reaches $7/9=77.8\%$ ($\geq75\%$, passing), comparable in magnitude to the CIFAR substitute's own flip accuracy ($83.3\%$ with all four mechanisms, Section~\ref{sec:baseline-comparison}; this replication uses the original three, since it predates quantization being added to the main comparison). The full predictor comparison recovers the Section~\ref{sec:baseline-comparison} pattern as well: the three direction-aware predictors cluster together (full $\rho=0.747$, diagonal $\rho=0.728$, first-order-only $\rho=0.740$; all $80.6\%$ sign-agreement, all $7/9$ flip accuracy), while the two magnitude-only predictors are markedly worse and the isotropic correlation is not statistically significant (isotropic $\rho=0.143$, $p=0.40$; norm bound and isotropic both $4/9=44.4\%$ flip accuracy). We find the same qualitative ordering as Table~\ref{tab:flips} and Figure~\ref{fig:baseline-comparison}, on a real dataset, under a single, independently-motivated change to the training recipe.

\begin{table}[h]
\centering
\caption{Waterbirds replication, both training regimes (Appendix~\ref{app:waterbirds}). Flip accuracy is directional-predictor (full-HVP) accuracy on cross-mechanism flip cases.}
\label{tab:waterbirds}
\begin{tabular}{@{}lrrrr@{}}
\toprule
Regime & Train acc. & Pooled $\rho$ & Flip seeds & Flip accuracy \\
\midrule
8 epochs (overfit) & $100\%$ & $0.559$ & $2/3$ & $50.0\%$ \\
3 epochs (early-stopped) & $99.3$--$99.5\%$ & $0.747$ & $3/3$ & $77.8\%$ \\
\bottomrule
\end{tabular}
\end{table}

\textbf{Interpretation.} This replication supports the paper's central claim, that compression mechanisms harm different groups by direction and not merely magnitude, on a real dataset, under the training regime where the theory's locality assumption is plausible. It also turns Section~\ref{sec:validity-frontier}'s open question from speculation into a concrete, if partial, answer: the boundary of validity is not solely a function of the compression severity $\|\delta\|/\|\theorig\|$, which we held fixed at $3\%$ across both regimes reported here. A model's own degree of overfitting at $\theorig$ (equivalently, how far a group's pre-compression loss already sits from the region where a local quadratic expansion is accurate) is at least as important, and the two overfitting regimes we test bracket a transition from the comparison failing outright (even the direction-vs-magnitude distinction disappears) to it replicating our main result closely. Neither run clears the pooled-Spearman threshold as cleanly as the CIFAR substitute does ($\rho=0.878$ on the original three mechanisms this replication is matched to, $\rho=0.901$ with quantization included; Section~\ref{sec:baseline-comparison}), which we attribute to Waterbirds' much larger, more heterogeneous ResNet-18-at-$128^2$ parameter space and smaller per-group evaluation samples (Section~\ref{sec:setup}) relative to the controlled substitute, rather than to a failure of the direction-vs-magnitude claim itself, which the flip-case accuracy (the harder, more diagnostic test; Figure~\ref{fig:baseline-comparison}'s caption) replicates closely. We did not search over training recipes beyond this one ablation; a systematic sweep over degree of overfitting, in the spirit of Section~\ref{sec:validity-frontier}'s proposed severity sweep, is left to future work.

\subsection{Quantization as the isotropic special case}
\label{app:quantization}

Section~\ref{sec:isotropic-special-case} predicts that post-training quantization, having an approximately isotropic perturbation, should be better described by the isotropic account than pruning or low-rank truncation are. We add quantization (uniform scalar rounding to a single global step size, Section~\ref{sec:setup}) as a fourth mechanism to the main comparison and test this.

\textbf{Quantization strengthens the headline numbers.} Calibrated to the same $\|\delta\|/\|\theorig\|\approx3\%$ target as the other three mechanisms (Table~\ref{tab:calibration}), quantization produces a flip in the most-harmed group relative to at least one other mechanism in every seed (Table~\ref{tab:flips}), and its predicted-vs-actual values from the full directional predictor are the tightest of any mechanism we test: e.g.\ seed 0, group cat+blue, actual $+0.0143$ vs.\ predicted $+0.0146$; seed 1, group dog+red, actual $-0.0208$ vs.\ predicted $-0.0222$; seed 2, group dog+red, actual $-0.0156$ vs.\ predicted $-0.0170$. Adding it raises the pooled Spearman correlation for all three direction-aware predictors from $\rho\approx0.88$ (three mechanisms, Section~\ref{sec:baseline-comparison}) to $\rho\approx0.90$ ($n=48$) and flip-case accuracy from $77.8\%$ to $83.3\%$ ($10/12$); the magnitude-only predictors remain at chance ($25.0\%$, $3/12$).

\textbf{The more specific isotropic-curvature prediction is a partial confirmation.} The test above uses the \emph{full} isotropic predictor, $\|g_a\|\|\delta\| + \tfrac12\mathrm{tr}(H_a)\|\delta\|^2/d$ (Section~\ref{sec:setup}), which is dominated by its sign-blind first-order term regardless of mechanism (it is at chance for all four mechanisms individually, not only pooled), so it cannot by itself confirm or refute a claim specifically about the \emph{curvature} term. We isolate that term directly: for each (seed, group, mechanism) we compare the true quadratic form $\delta^\top H_a\delta$ (recovered exactly from the full and first-order predictors already computed, since $\delta^\top H_a\delta = 2(\text{pred}_{\text{full}} - \text{pred}_{\text{first}})$) against its isotropic approximation $\mathrm{tr}(H_a)\|\delta\|^2/d$, per mechanism, across all $12$ (seed, group) pairs:

\begin{center}
\begin{tabular}{@{}lrr@{}}
\toprule
Mechanism & Correlation(true, isotropic) & Median $|\text{ratio}-1|$ \\
\midrule
Unstructured & $-0.01$ & $1.57$ \\
Channel & $-0.25$ & $1.00$ \\
Low-rank (SVD) & $0.66$ & $0.95$ \\
Quantization & $0.29$ & $\mathbf{0.65}$ \\
\bottomrule
\end{tabular}
\end{center}

Quantization has the lowest typical relative error between the true and isotropic-approximated curvature term (median ratio within $65\%$ of $1$, versus $95$--$157\%$ for the other three), the direction Section~\ref{sec:isotropic-special-case} predicts. But its correlation across the $12$ pairs ($0.29$) does not cleanly exceed every other mechanism's: low-rank truncation's is higher ($0.66$). We do not read this as refuting the isotropic-special-case argument, but we do not claim it as a clean confirmation either. With only $12$ points per mechanism, spanning three seeds whose Hessian traces differ by roughly an order of magnitude (Table~\ref{tab:calibration}'s seed-to-seed severities hint at the same underlying variation), a single correlation estimate is noisy, and the two metrics disagree on which non-quantization mechanism is the closest runner-up. We report both metrics rather than the more favorable one alone, and leave a larger-sample version of this specific test (more seeds, or evaluating within each seed separately rather than pooling across seeds of very different curvature scale) to future work.

\subsection{Accuracy effect sizes at the most-harmed group}
\label{app:accuracy-effect-sizes}

The paper reports loss throughout, since loss is the continuous quantity $\widehat{R}(a)$ (Section~\ref{sec:theory}) predicts directly and accuracy is not. Table~\ref{tab:accuracy-effects} reports, for each (seed, mechanism) pair in the main comparison, the accuracy of the actual most-harmed group (Table~\ref{tab:flips}) before and after compression.

\begin{table}[h]
\centering
\caption{Accuracy of the actual most-harmed group, before and after compression, for each (seed, mechanism) pair. $\Delta$acc is compressed minus base accuracy.}
\label{tab:accuracy-effects}
\begin{tabular}{@{}llrrr@{}}
\toprule
Seed & Mechanism & Base acc. & Compressed acc. & $\Delta$acc \\
\midrule
\multirow{4}{*}{0} & Unstructured & $0.000$ & $0.000$ & $+0.000$ \\
 & Channel & $0.000$ & $0.000$ & $+0.000$ \\
 & Low-rank (SVD) & $0.000$ & $0.000$ & $+0.000$ \\
 & Quantization & $0.000$ & $0.000$ & $+0.000$ \\
\addlinespace
\multirow{4}{*}{1} & Unstructured & $0.127$ & $0.127$ & $+0.000$ \\
 & Channel & $0.988$ & $0.985$ & $-0.004$ \\
 & Low-rank (SVD) & $0.127$ & $0.127$ & $+0.000$ \\
 & Quantization & $0.127$ & $0.127$ & $+0.000$ \\
\addlinespace
\multirow{4}{*}{2} & Unstructured & $0.113$ & $0.120$ & $+0.007$ \\
 & Channel & $0.113$ & $0.067$ & $-0.047$ \\
 & Low-rank (SVD) & $0.988$ & $0.986$ & $-0.002$ \\
 & Quantization & $0.113$ & $0.107$ & $-0.007$ \\
\bottomrule
\end{tabular}
\end{table}

Most entries show little or no accuracy change at the most-harmed group, including several seed-0 rows where the group's accuracy is already $0\%$ before compression. This is the floor effect Section~\ref{sec:setup} already describes: at the $85\%$ correlation strength we use, several minority groups' pre-compression accuracy is already at or near the benchmark's floor, so there is often little accuracy left to lose even as the group's loss, the quantity $\widehat{R}(a)$ actually predicts, continues to increase measurably (a model can become more confidently wrong, raising cross-entropy loss, while its argmax prediction, and hence accuracy, does not change). Where pre-compression accuracy is not already near floor (seed 1 channel pruning, seed 2 channel pruning and low-rank truncation), the accuracy drops are modest ($0.4$ to $4.7$ percentage points) but not zero, and the largest of them (seed 2, channel pruning, $-4.7$ points) occurs at exactly the (seed, mechanism) pair with the single largest actual loss increase among all $12$ in Table~\ref{tab:flips}. We read this as consistent with, not contradicting, the paper's central claim, but it is a real scope limitation: in a benchmark where minority-group accuracy sits further from floor, the second-order theory's loss predictions might translate more uniformly into accuracy effect sizes, and we have not tested that regime directly.

\subsection{Severity sweep: the validity frontier}
\label{app:severity-sweep}

Section~\ref{sec:validity-frontier} flags, as an open question, where across a range of $\|\delta\|/\|\theorig\|$ the directional predictor's accuracy degrades, beyond the single $3\%$ target used throughout the main comparison. We ran this sweep directly: holding the three CIFAR-substitute checkpoints fixed (no retraining, ruling out the overfitting confound Appendix~\ref{app:waterbirds} identifies as a separate axis), we recalibrated all four mechanisms to four severity targets, $\|\delta\|/\|\theorig\| \in \{1\%, 3\%, 10\%, 25\%\}$ (achieved within $0.01$ percentage points of each target in every case), and recomputed the full five-predictor comparison at each level. To keep the sweep's total compute tractable we used $50$ Hutchinson probes for the diagonal and isotropic predictors at each level, rather than the main comparison's $100$; Section~\ref{sec:setup}'s convergence check already shows this halving does not materially change the comparison's conclusions.

\begin{table}[h]
\centering
\caption{Severity sweep: pooled Spearman correlation and flip-case accuracy at four relative-perturbation targets ($n=48$ per level, $12$ flip cases per level).}
\label{tab:severity-sweep}
\begin{tabular}{@{}rrrrrrr@{}}
\toprule
Target $\|\delta\|/\|\theorig\|$ & Achieved & $\rho$(full) & $\rho$(first) & $\rho$(diag) & $\rho$(iso) & Acc.\ (full / iso) \\
\midrule
$1\%$  & $1.02\%$  & $0.997$ & $0.997$ & $0.970$ & $0.096$  & $100.0\%$ / $50.0\%$ \\
$3\%$  & $2.99\%$  & $0.901$ & $0.902$ & $0.893$ & $-0.038$ & $83.3\%$ / $25.0\%$ \\
$10\%$ & $10.00\%$ & $0.973$ & $0.935$ & $0.932$ & $-0.251$ & $83.3\%$ / $25.0\%$ \\
$25\%$ & $25.00\%$ & $0.976$ & $0.860$ & $0.879$ & $-0.337$ & $58.3\%$ / $25.0\%$ \\
\bottomrule
\end{tabular}
\end{table}

The $3\%$ row reproduces the main comparison's numbers closely (Section~\ref{sec:baseline-comparison}: $\rho=0.901$ for the full predictor, $83.3\%$ flip accuracy), confirming the sweep's independently-run calibration and predictor pipeline agree with the main result where they overlap.

Two findings emerge, and they do not tell the same story. First, the full directional predictor's \emph{pooled correlation} with actual loss does not monotonically degrade with severity: it is highest at the mildest severity we test ($\rho=0.997$ at $1\%$), dips at $3\%$, then recovers to $\rho\approx0.97$--$0.98$ at both $10\%$ and $25\%$. Taken alone, pooled correlation would suggest the second-order theory remains highly predictive even at $25\%$ relative perturbation, a severity roughly eight times larger than our main target. Second, \emph{flip-case accuracy}, the harder test that asks the predictor to correctly name the single most-harmed group rather than merely track loss changes in the right order, tells a different story: it falls from a perfect $100\%$ at $1\%$ severity to $58.3\%$ at $25\%$, approaching but not reaching the $25\%$ chance level. We read this divergence as itself informative about what "validity" means for this theory: the second-order expansion continues to rank-order groups' loss changes well even at large $\|\delta\|$, but its ability to resolve which single group is \emph{most} harmed, the practically relevant question, degrades well before the pooled correlation does. A practitioner relying on pooled correlation alone could be misled into trusting the predictor's discrete recommendation at a severity where it is no longer reliable.

The isotropic predictor's behavior is also informative in its own right: rather than remaining merely uninformative as severity grows, it becomes progressively more \emph{anti}-correlated with actual loss ($\rho=0.096$ at $1\%$, falling to $\rho=-0.337$ at $25\%$), while its flip-case accuracy stays pinned at the mechanism-independent norm bound's chance level ($25.0\%$) throughout. This is consistent with Section~\ref{sec:setup}'s observation that the isotropic predictor's sign-blind first-order term dominates its behavior regardless of mechanism: at larger $\|\delta\|$ that term's magnitude grows faster than the trace-based curvature term it is added to, making the predictor's ranking progressively less related to, and eventually inversely related to, the groups' actual loss ordering.

We tested four severity levels per seed on the original three checkpoints (CIFAR substitute); as in the main comparison, our evidence base is $3$ seeds, and a replication across more seeds, more severity levels (particularly between $3\%$ and $10\%$, where the two metrics' divergence seems to begin), or on Waterbirds would strengthen confidence in exactly where flip-case accuracy crosses below the paper's $75\%$ gate threshold.

\subsection{Disjoint predictor/evaluation split}
\label{app:disjoint-split}

Section~\ref{sec:setup} notes that the five predictors are computed on a $50$-example batch per group, which partially overlaps with the larger sample used to measure each group's actual loss before and after compression ($850$/$850$/$150$/$150$ examples for the four groups; $5.9\%$ overlap for majority groups, $33.3\%$ for minority groups). To rule out shared sampling noise as an explanation for the main result, we reran the comparison with the two uses of the test set made fully disjoint: for every group, examples $[0,50)$ are used only to compute the predictors, and a different, non-overlapping set of $50$ examples, $[50,100)$, are used only to measure actual loss before and after compression. Severities are not recalibrated; we reuse the exact severities already found for Table~\ref{tab:calibration}, so this isolates the effect of the overlap alone.

The result is nearly identical to the original (overlapping) comparison. Pooled Spearman: full directional $\rho=0.875$ ($p=4.6\times10^{-16}$, vs.\ $\rho=0.901$ originally), diagonal $\rho=0.890$ (vs.\ $0.901$), first-order-only $\rho=0.877$ (vs.\ $0.904$), isotropic $\rho=-0.067$ (vs.\ $-0.041$, still not significant, $p=0.65$). Flip-case accuracy is unchanged for the three direction-aware predictors ($10/12=83.3\%$, identical to the original) and for the isotropic predictor ($3/12=25.0\%$, chance). Separately, we measured how much the actual-loss values themselves change between the two halves: the Spearman correlation between each group's full-test-set-average actual loss (used in the main comparison) and its loss computed on just the disjoint second half is $\rho=0.985$ ($p=5.5\times10^{-37}$, $n=48$), indicating the actual-loss measurement itself is highly stable across different subsamples of the same group, independent of which half the predictor also used. We conclude the predictor/evaluation overlap flagged in Section~\ref{sec:setup} is not driving the main result.

\subsection{Weight-decay ablation}
\label{app:weight-decay}

Proposition~\ref{prop:stationarity} assumes $\theorig$ is a stationary point of the \emph{unregularized} training loss; our main experiments satisfy this by using no weight decay. To test whether the paper's empirical findings, as opposed to the proposition's analytic guarantee, depend on this choice, we retrained all three seeds with standard SGD weight decay ($5\times10^{-4}$), otherwise following Section~\ref{sec:setup}'s procedure unchanged, and reran the full gate and five-predictor comparison.

The gate passes cleanly on all three criteria: pooled Spearman $\rho=0.975$ ($p<10^{-30}$, $n=48$, exceeding the $0.80$ threshold), all $3$ seeds produce a flip in the most-harmed group across mechanisms (the same $3/3$ rate as the main no-weight-decay result), and directional-predictor accuracy on these $12$ flip cases is $9/12=75.0\%$ (meeting the $75\%$ threshold, though below the main result's $83.3\%$). The full five-predictor comparison recovers the same qualitative pattern: the three direction-aware predictors cluster together and tie at $75.0\%$ flip-case accuracy (full $\rho=0.975$, diagonal $\rho=0.981$, first-order-only $\rho=0.974$; all $95.8\%$ sign-agreement), while the two magnitude-only predictors remain at chance (isotropic $\rho=-0.163$, $p=0.27$, not statistically significant; both isotropic and the norm bound at $25.0\%$, $3/12$). If anything, the direction-aware predictors' correlations with actual loss are slightly \emph{higher} under weight decay than without it ($\rho\approx0.97$--$0.98$ vs.\ $\rho\approx0.88$--$0.90$ in Section~\ref{sec:baseline-comparison}), though flip-case accuracy is slightly lower ($75.0\%$ vs.\ $83.3\%$).

We read this as a clean, if not exhaustive, confirmation that the paper's central empirical findings, direction over magnitude, and no measurable benefit from curvature beyond the first-order term, do not depend on training without weight decay, even though Proposition~\ref{prop:stationarity}'s analytic guarantee (gradients summing to zero at $\theorig$) technically does not hold exactly once weight decay is present. We tested a single standard weight-decay value ($5\times10^{-4}$); a sweep across weight-decay strengths, in the spirit of Section~\ref{sec:validity-frontier}'s proposed severity sweep, is left to future work.

\subsection{Full derivations}
\label{app:derivations}

\paragraph{Proof of Proposition~\ref{prop:stationarity}.}
Let $\theorig$ be a stationary point of the unregularized, population-weighted training loss $L(\theta) = \sum_a w_a L_a(\theta)$, so $\nabla L(\theorig) = \sum_a w_a g_a = 0$, where $g_a = \nabla_\theta L_a(\theorig)$ and $w_a > 0$ for every group $a$, $\sum_a w_a = 1$.

\emph{First-order cancellation.} For any perturbation $\delta \in \RR^d$, taking the inner product of the stationarity condition with $\delta$ gives
\[
\sum_a w_a \big(g_a^\top \delta\big) \;=\; \Big(\sum_a w_a g_a\Big)^{\!\top} \delta \;=\; 0 .
\]
This is the weighted sum of the groups' first-order excess-loss terms under $\delta$, and it vanishes identically, for \emph{every} $\delta$, purely as a consequence of $\theorig$ being stationary for the pooled objective.

\emph{Two groups.} With exactly two groups $a,b$ ($w_a,w_b>0$, $w_a+w_b=1$), the cancellation condition reads $w_a(g_a^\top\delta) + w_b(g_b^\top\delta) = 0$, i.e.\ $g_a^\top\delta = -(w_b/w_a)\,g_b^\top\delta$. Since $w_a,w_b>0$, this forces $g_a^\top\delta$ and $g_b^\top\delta$ to carry opposite sign whenever either is nonzero (both are zero only if $\delta$ is orthogonal to both gradients). So for any $\delta$ not orthogonal to both $g_a$ and $g_b$, one group's first-order excess loss is positive and the other's is negative.

\emph{More than two groups.} For $k>2$ groups, $\sum_a w_a(g_a^\top\delta)=0$ is a single linear constraint on the $k$ numbers $\{g_a^\top\delta\}_{a=1}^k$, weighted by the positive weights $\{w_a\}$. A weighted sum of terms with strictly positive weights can vanish only if the terms do not all share one strict sign: at least one must be $\le 0$ and at least one $\ge 0$, unless every term is exactly zero. Beyond this the constraint is silent about any individual group's sign: it is equally consistent with $k-1$ groups all having $g_a^\top\delta>0$, provided the remaining group's (weighted) term is negative enough to offset them. Stationarity thus guarantees \emph{some} sign split exists but does not determine \emph{which} groups land on which side; that is decided by $\delta$'s direction relative to each $g_a$, not by stationarity alone.

\emph{Why this makes the ranking mechanism-dependent.} Fix the (mechanism-independent) gradients $\{g_a\}$. For each group $a$, the hyperplane $\{\delta : g_a^\top\delta = 0\}$ splits $\RR^d$ into the half-space where $a$'s first-order excess loss is positive and the half-space where it is negative. Two perturbation directions $\delta_1,\delta_2$, produced for instance by two different compression mechanisms even at matched $\|\delta_1\|=\|\delta_2\|$, can fall on opposite sides of group $a$'s hyperplane while falling on the \emph{same} side of group $b$'s hyperplane, or vice versa, whenever $g_a$ and $g_b$ are not parallel, which is the generic case. This is precisely the mechanism that can flip which group is helped and which is harmed at first order, and it is fully compatible with stationarity: stationarity fixes only the single weighted-sum constraint above, leaving the arrangement of these $k$ hyperplanes, and therefore which side of each one a given mechanism's $\delta$ falls on, entirely open. This is the formal sense in which Proposition~\ref{prop:stationarity} \emph{constrains} direction without \emph{fixing} it, and it is the reason the empirical question of Section~\ref{sec:results}, which group a given mechanism harms most, is not already settled by the stationarity of $\theorig$ alone. $\blacksquare$

\paragraph{Full Gauss--Newton derivation (Section~\ref{sec:gauss-newton}), multi-class case.}
Let $z\in\RR^C$ be the pre-softmax logits for a $C$-class example with true label $y$, $p=\mathrm{softmax}(z)$, and $\ell(z) = -\log p_y$ the cross-entropy loss; write $e_y$ for the one-hot indicator of $y$.

\emph{Gradient and Hessian in logit space.} Since $\ell(z) = -z_y + \log\sum_k e^{z_k}$,
\[
\frac{\partial \ell}{\partial z_j} = p_j - e_y[j], \qquad \frac{\partial^2\ell}{\partial z_j\partial z_l} = p_j\big(\mathbb{1}[j=l] - p_l\big),
\]
so $\nabla_z\ell = p-e_y$ and $\nabla^2_z\ell = \mathrm{diag}(p) - pp^\top$, the covariance matrix of a one-hot draw from $\mathrm{Categorical}(p)$ and hence positive semi-definite.

\emph{Chain rule into parameter space.} Let $z(\theta)$ be the network's logit map and $J = \partial z/\partial\theta \in \RR^{C\times d}$ its Jacobian at $\theorig$. By the chain rule,
\[
\frac{\partial \ell}{\partial\theta_m} = \sum_j \frac{\partial\ell}{\partial z_j}\frac{\partial z_j}{\partial\theta_m}, \qquad
\frac{\partial^2\ell}{\partial\theta_m\partial\theta_n} = \sum_{j,l}\frac{\partial^2\ell}{\partial z_j\partial z_l}\frac{\partial z_j}{\partial\theta_m}\frac{\partial z_l}{\partial\theta_n} \;+\; \sum_j \frac{\partial\ell}{\partial z_j}\frac{\partial^2 z_j}{\partial\theta_m\partial\theta_n},
\]
i.e., in matrix form the exact per-example Hessian decomposes as
\[
H \;=\; \underbrace{J^\top\big(\mathrm{diag}(p)-pp^\top\big)J}_{\text{Gauss--Newton term}} \;+\; \underbrace{\textstyle\sum_j (p_j - e_y[j])\,\nabla^2_\theta z_j}_{\text{residual term}} .
\]
The Gauss--Newton approximation keeps only the first term, $H_{\mathrm{GN}} = J^\top(\mathrm{diag}(p)-pp^\top)J$, which is PSD regardless of the network's nonlinearity; the residual term involves the second derivative of the (generally nonlinear) logit map itself and is weighted by the residual $(p-e_y)$, the same quantity that governs the gradient. The exact Hessian we use empirically (Section~\ref{sec:setup}) retains this residual term; the Gauss--Newton approximation does not.

\emph{Binary case.} With $C=2$ and a scalar logit $z$ for the true class (so $p$ is the predicted probability of the true label), $\partial\ell/\partial z = p-1$ and $\partial^2\ell/\partial z^2 = p(1-p)$, recovering $\|g_i\| = |p_i-1|\cdot\|J_i\|$ exactly and $H_{\mathrm{GN},i} = p_i(1-p_i)\|J_i\|^2$, as used in Section~\ref{sec:gauss-newton}.

\emph{Decoupling at any simplex vertex.} As $p\to e_k$ for \emph{any} class $k$ (not only $k=y$), $\mathrm{diag}(p)-pp^\top \to \mathrm{diag}(e_k) - e_ke_k^\top = 0$: the Gauss--Newton curvature vanishes at every vertex of the simplex, since a (near-)deterministic categorical distribution has (near-)zero variance. Meanwhile the gradient-space residual $\|p-e_y\|\to\|e_k-e_y\|=\sqrt{2}$ for every $k\ne y$, and does not vanish. So for $k\ne y$, a confidently \emph{misclassified} example, the Gauss--Newton curvature term collapses to zero while the gradient stays bounded away from zero: the same decoupling as the binary case ($p_i\to 0$), extended to every non-target vertex. This is also exactly the regime where the discarded residual term above is largest, since it is weighted by the same $(p-e_y)$ that stays large while $H_{\mathrm{GN}}$ vanishes. This is the reason Section~\ref{sec:gauss-newton} treats this as motivation for where the isotropic approximation of Section~\ref{sec:isotropic-special-case} may be least trustworthy, rather than as a proof about the exact-Hessian predictor, which retains the residual term the Gauss--Newton approximation drops. $\blacksquare$

\end{document}
